% =============================================================================
%  ulms-preamble.tex — preamble ที่ใช้ร่วมกันทุกเอกสารของโครงการ ULMs
%
%  ใช้โดย:  main.tex (ข้อเสนอโครงการ) · progress.tex (รายงานความก้าวหน้า)
%
%  เอกสารที่เรียกใช้ต้องทำสองอย่างหลัง % =============================================================================
%  ulms-preamble.tex — preamble ที่ใช้ร่วมกันทุกเอกสารของโครงการ ULMs
%
%  ใช้โดย:  main.tex (ข้อเสนอโครงการ) · progress.tex (รายงานความก้าวหน้า)
%
%  เอกสารที่เรียกใช้ต้องทำสองอย่างหลัง % =============================================================================
%  ulms-preamble.tex — preamble ที่ใช้ร่วมกันทุกเอกสารของโครงการ ULMs
%
%  ใช้โดย:  main.tex (ข้อเสนอโครงการ) · progress.tex (รายงานความก้าวหน้า)
%
%  เอกสารที่เรียกใช้ต้องทำสองอย่างหลัง % =============================================================================
%  ulms-preamble.tex — preamble ที่ใช้ร่วมกันทุกเอกสารของโครงการ ULMs
%
%  ใช้โดย:  main.tex (ข้อเสนอโครงการ) · progress.tex (รายงานความก้าวหน้า)
%
%  เอกสารที่เรียกใช้ต้องทำสองอย่างหลัง \input{ulms-preamble}:
%    1) \hypersetup{pdftitle={...}}  — ชื่อเอกสารใน metadata ของ PDF
%    2) \begin{document}
%  (pdftitle ถูกย้ายออกจากไฟล์นี้ เพราะเป็นค่าเฉพาะของแต่ละเอกสาร)
%
%  แก้ที่นี่ที่เดียว = มีผลกับทุกเอกสาร ห้ามก๊อป preamble ไปวางซ้ำ
% =============================================================================
% =============================================================================
%  main.tex  —  เอกสารข้อเสนอโครงการ (Academic Project Proposal)
%  ระบบบริหารจัดการการยืม–คืนอุปกรณ์มหาวิทยาลัย (ULMs)
%
%  IMPORTANT: This document contains Thai text and MUST be compiled with
%  XeLaTeX (not pdflatex). Thai needs a Unicode font + ICU line-breaking.
%      Build:  latexmk main.tex        (uses .latexmkrc -> xelatex)
%      or:     xelatex main.tex        (run twice for the table of contents)
% =============================================================================

\documentclass[11pt, a4paper]{article}

% ---- Fonts & Thai support ----------------------------------------------------
% fontspec lets XeLaTeX use any system font. Laksaman is a complete serif face
% that covers BOTH Thai and Latin glyphs, so one \setmainfont handles the whole
% (Thai + English) document without font switching.
\usepackage{fontspec}

% Thai has no spaces between words, so TeX cannot find line-break points on its
% own. These two XeTeX primitives switch on ICU's dictionary-based Thai word
% segmentation, which is what makes Thai paragraphs wrap correctly.
\XeTeXlinebreaklocale "th"
\XeTeXlinebreakskip = 0pt plus 0.1pt

\setmainfont{Laksaman}

% ---- Page geometry & paragraph style ----------------------------------------
\usepackage[margin=2.5cm]{geometry}
\usepackage{parskip}          % block paragraphs (no indent, gap between) — common
                              % and more readable for Thai documents
\linespread{1.15}             % a little extra leading; Thai has tall tone marks

% ---- Colours (monochrome / grayscale palette) ------------------------------
\usepackage[table]{xcolor}    % 'table' option enables \rowcolor for table headers
% Monochrome (grayscale) palette. Names kept as-is so every reference still
% works; only the tones changed from the original green/blue accents.
\definecolor{kugreen}{HTML}{1A1A1A}   % near-black — section headings & title
\definecolor{kudark}{HTML}{333333}    % dark gray — subsection/subsubsection titles
\definecolor{headbg}{HTML}{EAEAEA}    % light gray background for table header rows
\definecolor{linkblue}{HTML}{404040}  % medium-dark gray — links / citations

% ---- Section heading styling -------------------------------------------------
% 'nobottomtitles*' stops a heading from being printed near the bottom of a
% page: if less than \bottomtitlespace of room is left, the heading (and its
% following text) moves to the next page. This is titlesec's built-in, robust
% replacement for the earlier \needspace hack — the latter injected vertical
% glue that clashed with longtable page-splitting ("Infinite glue shrinkage").
\usepackage[nobottomtitles*]{titlesec}
\titleformat{\section}
  {\color{kugreen}\Large\bfseries}{\thesection}{0.6em}{}
\titleformat{\subsection}
  {\color{kudark}\large\bfseries}{\thesubsection}{0.5em}{}
\titleformat{\subsubsection}
  {\color{kudark}\normalsize\bfseries}{\thesubsubsection}{0.5em}{}
\titlespacing*{\section}{0pt}{1.4\baselineskip}{0.5\baselineskip}

% Minimum room a heading needs at the bottom of a page; if only ~1–3 lines are
% left (less than this), the heading is bumped to the next page instead of being
% stranded above its table/text.
\renewcommand{\bottomtitlespace}{4\baselineskip}

% ---- Hierarchical indentation by heading level ------------------------------
% Each heading level shifts BOTH the heading and the body text under it further
% to the right: section = 0, subsection (x.y) = 1 step, subsubsection (x.y.z) =
% 2 steps. We do this by having the sectioning commands set \leftskip (a global
% left margin that carries through every following line until the next heading
% resets it). \par first, so the *previous* paragraph is not re-indented.
% \leftskip indents plain paragraphs; \@totalleftmargin makes list environments
% (itemize/enumerate) start from the same indent instead of the page margin.
\makeatletter
\newlength{\lvlindent}
\setlength{\lvlindent}{1.8em}          % one indent step
\let\ulmsSecIndent\section
\renewcommand{\section}{\par\global\leftskip=0pt\global\@totalleftmargin=0pt\relax\ulmsSecIndent}
\let\ulmsSubsecIndent\subsection
\renewcommand{\subsection}{\par\global\leftskip=\lvlindent\global\@totalleftmargin=\lvlindent\relax\ulmsSubsecIndent}
\let\ulmsSubsubsecIndent\subsubsection
\renewcommand{\subsubsection}{\par\global\leftskip=2\lvlindent\global\@totalleftmargin=2\lvlindent\relax\ulmsSubsubsecIndent}
\makeatother

% Indent the headings themselves to match (titlesec ignores \leftskip for the
% heading line, so its left margin is set here via the first \titlespacing arg).
\titlespacing*{\subsection}   {\lvlindent} {1.0\baselineskip}{0.4\baselineskip}
\titlespacing*{\subsubsection}{2\lvlindent}{0.8\baselineskip}{0.3\baselineskip}

% ---- Lists (tighter, custom bullets) ----------------------------------------
\usepackage{enumitem}
\setlist[itemize]{leftmargin=1.6em, itemsep=2pt, topsep=3pt}
\setlist[enumerate]{leftmargin=1.9em, itemsep=2pt, topsep=3pt}

% ---- Tables ------------------------------------------------------------------
\usepackage{array}
\usepackage{booktabs}         % \toprule / \midrule / \bottomrule
\usepackage{tabularx}         % auto-width columns that wrap long Thai text
\usepackage{longtable}        % tables that can break across pages
\usepackage{multirow}         % merged cells for the team-structure table

% Left-aligned wrapping column types:
%   L{width} — fixed-width ragged-right paragraph column
%   Y        — flexible ragged-right X column (for tabularx)
\newcolumntype{L}[1]{>{\raggedright\arraybackslash}p{#1}}
\newcolumntype{Y}{>{\raggedright\arraybackslash}X}
\renewcommand{\arraystretch}{1.3}

% Helpers for a styled header row inside tables.
%   \hrow  — starts a header row and shades it (must be FIRST in the row;
%            colortbl allows only one \rowcolor per row).
%   \thead — bold header-cell text.
\newcommand{\hrow}{\rowcolor{headbg}}
\newcommand{\thead}[1]{\textbf{#1}}

% \begin{keeptable} … \end{keeptable} — keeps a table and the bold label above
% it on the SAME page (used by the data dictionary in §5). A tabularx cannot
% split across pages, so LaTeX is free to break between the label and its
% table, stranding the label alone at the bottom of a page; wrapping the pair
% in a minipage makes them one unbreakable block.
%   * \leftskip is zeroed inside: the minipage box is already placed at the
%     current heading indent, so leaving it set would indent the content twice
%     and push the table past the right margin.
%   * \parskip is restored because minipage's \@parboxrestore zeroes it, which
%     would glue the table straight onto the label line.
% Tables inside must use \linewidth (= the minipage width), not \textwidth.
\newenvironment{keeptable}
  {\par\medskip\noindent
   \begin{minipage}{\dimexpr\textwidth-\leftskip\relax}%
   \leftskip=0pt\relax
   \setlength{\parskip}{0.5\baselineskip plus 2pt}}
  {\end{minipage}\par\medskip}

% ---- Table of contents styling ----------------------------------------------
% By default LaTeX gives dotted leaders to subsections but NOT to sections
% (they print bold with just a right-aligned page number). tocloft lets us turn
% on dots for the section level too, so every TOC line runs "….." to its page.
\usepackage{tocloft}
\renewcommand{\cftsecdotsep}{\cftdotsep}   % dotted leaders for section entries
\renewcommand{\cftsecleader}{\cftdotfill{\cftsecdotsep}}
\renewcommand{\cftsecpagefont}{\normalfont}  % keep page numbers un-bold

% ---- Figures -----------------------------------------------------------------
\usepackage{graphicx}
\graphicspath{{figures/}}

% Landscape (rotated) full-page floats — needed by the ER diagram in §5, which
% is far wider than it is tall and would be unreadable scaled to \textwidth.
\usepackage{rotating}

% ---- Flowcharts / diagrams (TikZ) -------------------------------------------
% Used for the borrow–return workflow diagram (§5). Libraries: geometric shapes
% for decision diamonds, arrows.meta for arrowheads, positioning for relative
% node placement, calc/fit/backgrounds for routing loops and grouping.
\usepackage{tikz}
\usetikzlibrary{shapes.geometric, arrows.meta, positioning, calc, fit, backgrounds}

% ---- Links & bookmarks -------------------------------------------------------
\usepackage{hyperref}
\hypersetup{
  colorlinks=true, linkcolor=black, citecolor=linkblue, urlcolor=linkblue,
  pdflang={th},
}

% Thai labels for auto-generated headings.
\renewcommand{\contentsname}{สารบัญ (Table of Contents)}
\renewcommand{\tablename}{ตาราง}
\renewcommand{\figurename}{รูปที่}

:
%    1) \hypersetup{pdftitle={...}}  — ชื่อเอกสารใน metadata ของ PDF
%    2) \begin{document}
%  (pdftitle ถูกย้ายออกจากไฟล์นี้ เพราะเป็นค่าเฉพาะของแต่ละเอกสาร)
%
%  แก้ที่นี่ที่เดียว = มีผลกับทุกเอกสาร ห้ามก๊อป preamble ไปวางซ้ำ
% =============================================================================
% =============================================================================
%  main.tex  —  เอกสารข้อเสนอโครงการ (Academic Project Proposal)
%  ระบบบริหารจัดการการยืม–คืนอุปกรณ์มหาวิทยาลัย (ULMs)
%
%  IMPORTANT: This document contains Thai text and MUST be compiled with
%  XeLaTeX (not pdflatex). Thai needs a Unicode font + ICU line-breaking.
%      Build:  latexmk main.tex        (uses .latexmkrc -> xelatex)
%      or:     xelatex main.tex        (run twice for the table of contents)
% =============================================================================

\documentclass[11pt, a4paper]{article}

% ---- Fonts & Thai support ----------------------------------------------------
% fontspec lets XeLaTeX use any system font. Laksaman is a complete serif face
% that covers BOTH Thai and Latin glyphs, so one \setmainfont handles the whole
% (Thai + English) document without font switching.
\usepackage{fontspec}

% Thai has no spaces between words, so TeX cannot find line-break points on its
% own. These two XeTeX primitives switch on ICU's dictionary-based Thai word
% segmentation, which is what makes Thai paragraphs wrap correctly.
\XeTeXlinebreaklocale "th"
\XeTeXlinebreakskip = 0pt plus 0.1pt

\setmainfont{Laksaman}

% ---- Page geometry & paragraph style ----------------------------------------
\usepackage[margin=2.5cm]{geometry}
\usepackage{parskip}          % block paragraphs (no indent, gap between) — common
                              % and more readable for Thai documents
\linespread{1.15}             % a little extra leading; Thai has tall tone marks

% ---- Colours (monochrome / grayscale palette) ------------------------------
\usepackage[table]{xcolor}    % 'table' option enables \rowcolor for table headers
% Monochrome (grayscale) palette. Names kept as-is so every reference still
% works; only the tones changed from the original green/blue accents.
\definecolor{kugreen}{HTML}{1A1A1A}   % near-black — section headings & title
\definecolor{kudark}{HTML}{333333}    % dark gray — subsection/subsubsection titles
\definecolor{headbg}{HTML}{EAEAEA}    % light gray background for table header rows
\definecolor{linkblue}{HTML}{404040}  % medium-dark gray — links / citations

% ---- Section heading styling -------------------------------------------------
% 'nobottomtitles*' stops a heading from being printed near the bottom of a
% page: if less than \bottomtitlespace of room is left, the heading (and its
% following text) moves to the next page. This is titlesec's built-in, robust
% replacement for the earlier \needspace hack — the latter injected vertical
% glue that clashed with longtable page-splitting ("Infinite glue shrinkage").
\usepackage[nobottomtitles*]{titlesec}
\titleformat{\section}
  {\color{kugreen}\Large\bfseries}{\thesection}{0.6em}{}
\titleformat{\subsection}
  {\color{kudark}\large\bfseries}{\thesubsection}{0.5em}{}
\titleformat{\subsubsection}
  {\color{kudark}\normalsize\bfseries}{\thesubsubsection}{0.5em}{}
\titlespacing*{\section}{0pt}{1.4\baselineskip}{0.5\baselineskip}

% Minimum room a heading needs at the bottom of a page; if only ~1–3 lines are
% left (less than this), the heading is bumped to the next page instead of being
% stranded above its table/text.
\renewcommand{\bottomtitlespace}{4\baselineskip}

% ---- Hierarchical indentation by heading level ------------------------------
% Each heading level shifts BOTH the heading and the body text under it further
% to the right: section = 0, subsection (x.y) = 1 step, subsubsection (x.y.z) =
% 2 steps. We do this by having the sectioning commands set \leftskip (a global
% left margin that carries through every following line until the next heading
% resets it). \par first, so the *previous* paragraph is not re-indented.
% \leftskip indents plain paragraphs; \@totalleftmargin makes list environments
% (itemize/enumerate) start from the same indent instead of the page margin.
\makeatletter
\newlength{\lvlindent}
\setlength{\lvlindent}{1.8em}          % one indent step
\let\ulmsSecIndent\section
\renewcommand{\section}{\par\global\leftskip=0pt\global\@totalleftmargin=0pt\relax\ulmsSecIndent}
\let\ulmsSubsecIndent\subsection
\renewcommand{\subsection}{\par\global\leftskip=\lvlindent\global\@totalleftmargin=\lvlindent\relax\ulmsSubsecIndent}
\let\ulmsSubsubsecIndent\subsubsection
\renewcommand{\subsubsection}{\par\global\leftskip=2\lvlindent\global\@totalleftmargin=2\lvlindent\relax\ulmsSubsubsecIndent}
\makeatother

% Indent the headings themselves to match (titlesec ignores \leftskip for the
% heading line, so its left margin is set here via the first \titlespacing arg).
\titlespacing*{\subsection}   {\lvlindent} {1.0\baselineskip}{0.4\baselineskip}
\titlespacing*{\subsubsection}{2\lvlindent}{0.8\baselineskip}{0.3\baselineskip}

% ---- Lists (tighter, custom bullets) ----------------------------------------
\usepackage{enumitem}
\setlist[itemize]{leftmargin=1.6em, itemsep=2pt, topsep=3pt}
\setlist[enumerate]{leftmargin=1.9em, itemsep=2pt, topsep=3pt}

% ---- Tables ------------------------------------------------------------------
\usepackage{array}
\usepackage{booktabs}         % \toprule / \midrule / \bottomrule
\usepackage{tabularx}         % auto-width columns that wrap long Thai text
\usepackage{longtable}        % tables that can break across pages
\usepackage{multirow}         % merged cells for the team-structure table

% Left-aligned wrapping column types:
%   L{width} — fixed-width ragged-right paragraph column
%   Y        — flexible ragged-right X column (for tabularx)
\newcolumntype{L}[1]{>{\raggedright\arraybackslash}p{#1}}
\newcolumntype{Y}{>{\raggedright\arraybackslash}X}
\renewcommand{\arraystretch}{1.3}

% Helpers for a styled header row inside tables.
%   \hrow  — starts a header row and shades it (must be FIRST in the row;
%            colortbl allows only one \rowcolor per row).
%   \thead — bold header-cell text.
\newcommand{\hrow}{\rowcolor{headbg}}
\newcommand{\thead}[1]{\textbf{#1}}

% \begin{keeptable} … \end{keeptable} — keeps a table and the bold label above
% it on the SAME page (used by the data dictionary in §5). A tabularx cannot
% split across pages, so LaTeX is free to break between the label and its
% table, stranding the label alone at the bottom of a page; wrapping the pair
% in a minipage makes them one unbreakable block.
%   * \leftskip is zeroed inside: the minipage box is already placed at the
%     current heading indent, so leaving it set would indent the content twice
%     and push the table past the right margin.
%   * \parskip is restored because minipage's \@parboxrestore zeroes it, which
%     would glue the table straight onto the label line.
% Tables inside must use \linewidth (= the minipage width), not \textwidth.
\newenvironment{keeptable}
  {\par\medskip\noindent
   \begin{minipage}{\dimexpr\textwidth-\leftskip\relax}%
   \leftskip=0pt\relax
   \setlength{\parskip}{0.5\baselineskip plus 2pt}}
  {\end{minipage}\par\medskip}

% ---- Table of contents styling ----------------------------------------------
% By default LaTeX gives dotted leaders to subsections but NOT to sections
% (they print bold with just a right-aligned page number). tocloft lets us turn
% on dots for the section level too, so every TOC line runs "….." to its page.
\usepackage{tocloft}
\renewcommand{\cftsecdotsep}{\cftdotsep}   % dotted leaders for section entries
\renewcommand{\cftsecleader}{\cftdotfill{\cftsecdotsep}}
\renewcommand{\cftsecpagefont}{\normalfont}  % keep page numbers un-bold

% ---- Figures -----------------------------------------------------------------
\usepackage{graphicx}
\graphicspath{{figures/}}

% Landscape (rotated) full-page floats — needed by the ER diagram in §5, which
% is far wider than it is tall and would be unreadable scaled to \textwidth.
\usepackage{rotating}

% ---- Flowcharts / diagrams (TikZ) -------------------------------------------
% Used for the borrow–return workflow diagram (§5). Libraries: geometric shapes
% for decision diamonds, arrows.meta for arrowheads, positioning for relative
% node placement, calc/fit/backgrounds for routing loops and grouping.
\usepackage{tikz}
\usetikzlibrary{shapes.geometric, arrows.meta, positioning, calc, fit, backgrounds}

% ---- Links & bookmarks -------------------------------------------------------
\usepackage{hyperref}
\hypersetup{
  colorlinks=true, linkcolor=black, citecolor=linkblue, urlcolor=linkblue,
  pdflang={th},
}

% Thai labels for auto-generated headings.
\renewcommand{\contentsname}{สารบัญ (Table of Contents)}
\renewcommand{\tablename}{ตาราง}
\renewcommand{\figurename}{รูปที่}

:
%    1) \hypersetup{pdftitle={...}}  — ชื่อเอกสารใน metadata ของ PDF
%    2) \begin{document}
%  (pdftitle ถูกย้ายออกจากไฟล์นี้ เพราะเป็นค่าเฉพาะของแต่ละเอกสาร)
%
%  แก้ที่นี่ที่เดียว = มีผลกับทุกเอกสาร ห้ามก๊อป preamble ไปวางซ้ำ
% =============================================================================
% =============================================================================
%  main.tex  —  เอกสารข้อเสนอโครงการ (Academic Project Proposal)
%  ระบบบริหารจัดการการยืม–คืนอุปกรณ์มหาวิทยาลัย (ULMs)
%
%  IMPORTANT: This document contains Thai text and MUST be compiled with
%  XeLaTeX (not pdflatex). Thai needs a Unicode font + ICU line-breaking.
%      Build:  latexmk main.tex        (uses .latexmkrc -> xelatex)
%      or:     xelatex main.tex        (run twice for the table of contents)
% =============================================================================

\documentclass[11pt, a4paper]{article}

% ---- Fonts & Thai support ----------------------------------------------------
% fontspec lets XeLaTeX use any system font. Laksaman is a complete serif face
% that covers BOTH Thai and Latin glyphs, so one \setmainfont handles the whole
% (Thai + English) document without font switching.
\usepackage{fontspec}

% Thai has no spaces between words, so TeX cannot find line-break points on its
% own. These two XeTeX primitives switch on ICU's dictionary-based Thai word
% segmentation, which is what makes Thai paragraphs wrap correctly.
\XeTeXlinebreaklocale "th"
\XeTeXlinebreakskip = 0pt plus 0.1pt

\setmainfont{Laksaman}

% ---- Page geometry & paragraph style ----------------------------------------
\usepackage[margin=2.5cm]{geometry}
\usepackage{parskip}          % block paragraphs (no indent, gap between) — common
                              % and more readable for Thai documents
\linespread{1.15}             % a little extra leading; Thai has tall tone marks

% ---- Colours (monochrome / grayscale palette) ------------------------------
\usepackage[table]{xcolor}    % 'table' option enables \rowcolor for table headers
% Monochrome (grayscale) palette. Names kept as-is so every reference still
% works; only the tones changed from the original green/blue accents.
\definecolor{kugreen}{HTML}{1A1A1A}   % near-black — section headings & title
\definecolor{kudark}{HTML}{333333}    % dark gray — subsection/subsubsection titles
\definecolor{headbg}{HTML}{EAEAEA}    % light gray background for table header rows
\definecolor{linkblue}{HTML}{404040}  % medium-dark gray — links / citations

% ---- Section heading styling -------------------------------------------------
% 'nobottomtitles*' stops a heading from being printed near the bottom of a
% page: if less than \bottomtitlespace of room is left, the heading (and its
% following text) moves to the next page. This is titlesec's built-in, robust
% replacement for the earlier \needspace hack — the latter injected vertical
% glue that clashed with longtable page-splitting ("Infinite glue shrinkage").
\usepackage[nobottomtitles*]{titlesec}
\titleformat{\section}
  {\color{kugreen}\Large\bfseries}{\thesection}{0.6em}{}
\titleformat{\subsection}
  {\color{kudark}\large\bfseries}{\thesubsection}{0.5em}{}
\titleformat{\subsubsection}
  {\color{kudark}\normalsize\bfseries}{\thesubsubsection}{0.5em}{}
\titlespacing*{\section}{0pt}{1.4\baselineskip}{0.5\baselineskip}

% Minimum room a heading needs at the bottom of a page; if only ~1–3 lines are
% left (less than this), the heading is bumped to the next page instead of being
% stranded above its table/text.
\renewcommand{\bottomtitlespace}{4\baselineskip}

% ---- Hierarchical indentation by heading level ------------------------------
% Each heading level shifts BOTH the heading and the body text under it further
% to the right: section = 0, subsection (x.y) = 1 step, subsubsection (x.y.z) =
% 2 steps. We do this by having the sectioning commands set \leftskip (a global
% left margin that carries through every following line until the next heading
% resets it). \par first, so the *previous* paragraph is not re-indented.
% \leftskip indents plain paragraphs; \@totalleftmargin makes list environments
% (itemize/enumerate) start from the same indent instead of the page margin.
\makeatletter
\newlength{\lvlindent}
\setlength{\lvlindent}{1.8em}          % one indent step
\let\ulmsSecIndent\section
\renewcommand{\section}{\par\global\leftskip=0pt\global\@totalleftmargin=0pt\relax\ulmsSecIndent}
\let\ulmsSubsecIndent\subsection
\renewcommand{\subsection}{\par\global\leftskip=\lvlindent\global\@totalleftmargin=\lvlindent\relax\ulmsSubsecIndent}
\let\ulmsSubsubsecIndent\subsubsection
\renewcommand{\subsubsection}{\par\global\leftskip=2\lvlindent\global\@totalleftmargin=2\lvlindent\relax\ulmsSubsubsecIndent}
\makeatother

% Indent the headings themselves to match (titlesec ignores \leftskip for the
% heading line, so its left margin is set here via the first \titlespacing arg).
\titlespacing*{\subsection}   {\lvlindent} {1.0\baselineskip}{0.4\baselineskip}
\titlespacing*{\subsubsection}{2\lvlindent}{0.8\baselineskip}{0.3\baselineskip}

% ---- Lists (tighter, custom bullets) ----------------------------------------
\usepackage{enumitem}
\setlist[itemize]{leftmargin=1.6em, itemsep=2pt, topsep=3pt}
\setlist[enumerate]{leftmargin=1.9em, itemsep=2pt, topsep=3pt}

% ---- Tables ------------------------------------------------------------------
\usepackage{array}
\usepackage{booktabs}         % \toprule / \midrule / \bottomrule
\usepackage{tabularx}         % auto-width columns that wrap long Thai text
\usepackage{longtable}        % tables that can break across pages
\usepackage{multirow}         % merged cells for the team-structure table

% Left-aligned wrapping column types:
%   L{width} — fixed-width ragged-right paragraph column
%   Y        — flexible ragged-right X column (for tabularx)
\newcolumntype{L}[1]{>{\raggedright\arraybackslash}p{#1}}
\newcolumntype{Y}{>{\raggedright\arraybackslash}X}
\renewcommand{\arraystretch}{1.3}

% Helpers for a styled header row inside tables.
%   \hrow  — starts a header row and shades it (must be FIRST in the row;
%            colortbl allows only one \rowcolor per row).
%   \thead — bold header-cell text.
\newcommand{\hrow}{\rowcolor{headbg}}
\newcommand{\thead}[1]{\textbf{#1}}

% \begin{keeptable} … \end{keeptable} — keeps a table and the bold label above
% it on the SAME page (used by the data dictionary in §5). A tabularx cannot
% split across pages, so LaTeX is free to break between the label and its
% table, stranding the label alone at the bottom of a page; wrapping the pair
% in a minipage makes them one unbreakable block.
%   * \leftskip is zeroed inside: the minipage box is already placed at the
%     current heading indent, so leaving it set would indent the content twice
%     and push the table past the right margin.
%   * \parskip is restored because minipage's \@parboxrestore zeroes it, which
%     would glue the table straight onto the label line.
% Tables inside must use \linewidth (= the minipage width), not \textwidth.
\newenvironment{keeptable}
  {\par\medskip\noindent
   \begin{minipage}{\dimexpr\textwidth-\leftskip\relax}%
   \leftskip=0pt\relax
   \setlength{\parskip}{0.5\baselineskip plus 2pt}}
  {\end{minipage}\par\medskip}

% ---- Table of contents styling ----------------------------------------------
% By default LaTeX gives dotted leaders to subsections but NOT to sections
% (they print bold with just a right-aligned page number). tocloft lets us turn
% on dots for the section level too, so every TOC line runs "….." to its page.
\usepackage{tocloft}
\renewcommand{\cftsecdotsep}{\cftdotsep}   % dotted leaders for section entries
\renewcommand{\cftsecleader}{\cftdotfill{\cftsecdotsep}}
\renewcommand{\cftsecpagefont}{\normalfont}  % keep page numbers un-bold

% ---- Figures -----------------------------------------------------------------
\usepackage{graphicx}
\graphicspath{{figures/}}

% Landscape (rotated) full-page floats — needed by the ER diagram in §5, which
% is far wider than it is tall and would be unreadable scaled to \textwidth.
\usepackage{rotating}

% ---- Flowcharts / diagrams (TikZ) -------------------------------------------
% Used for the borrow–return workflow diagram (§5). Libraries: geometric shapes
% for decision diamonds, arrows.meta for arrowheads, positioning for relative
% node placement, calc/fit/backgrounds for routing loops and grouping.
\usepackage{tikz}
\usetikzlibrary{shapes.geometric, arrows.meta, positioning, calc, fit, backgrounds}

% ---- Links & bookmarks -------------------------------------------------------
\usepackage{hyperref}
\hypersetup{
  colorlinks=true, linkcolor=black, citecolor=linkblue, urlcolor=linkblue,
  pdflang={th},
}

% Thai labels for auto-generated headings.
\renewcommand{\contentsname}{สารบัญ (Table of Contents)}
\renewcommand{\tablename}{ตาราง}
\renewcommand{\figurename}{รูปที่}

:
%    1) \hypersetup{pdftitle={...}}  — ชื่อเอกสารใน metadata ของ PDF
%    2) \begin{document}
%  (pdftitle ถูกย้ายออกจากไฟล์นี้ เพราะเป็นค่าเฉพาะของแต่ละเอกสาร)
%
%  แก้ที่นี่ที่เดียว = มีผลกับทุกเอกสาร ห้ามก๊อป preamble ไปวางซ้ำ
% =============================================================================
% =============================================================================
%  main.tex  —  เอกสารข้อเสนอโครงการ (Academic Project Proposal)
%  ระบบบริหารจัดการการยืม–คืนอุปกรณ์มหาวิทยาลัย (ULMs)
%
%  IMPORTANT: This document contains Thai text and MUST be compiled with
%  XeLaTeX (not pdflatex). Thai needs a Unicode font + ICU line-breaking.
%      Build:  latexmk main.tex        (uses .latexmkrc -> xelatex)
%      or:     xelatex main.tex        (run twice for the table of contents)
% =============================================================================

\documentclass[11pt, a4paper]{article}

% ---- Fonts & Thai support ----------------------------------------------------
% fontspec lets XeLaTeX use any system font. Laksaman is a complete serif face
% that covers BOTH Thai and Latin glyphs, so one \setmainfont handles the whole
% (Thai + English) document without font switching.
\usepackage{fontspec}

% Thai has no spaces between words, so TeX cannot find line-break points on its
% own. These two XeTeX primitives switch on ICU's dictionary-based Thai word
% segmentation, which is what makes Thai paragraphs wrap correctly.
\XeTeXlinebreaklocale "th"
\XeTeXlinebreakskip = 0pt plus 0.1pt

\setmainfont{Laksaman}

% ---- Page geometry & paragraph style ----------------------------------------
\usepackage[margin=2.5cm]{geometry}
\usepackage{parskip}          % block paragraphs (no indent, gap between) — common
                              % and more readable for Thai documents
\linespread{1.15}             % a little extra leading; Thai has tall tone marks

% ---- Colours (monochrome / grayscale palette) ------------------------------
\usepackage[table]{xcolor}    % 'table' option enables \rowcolor for table headers
% Monochrome (grayscale) palette. Names kept as-is so every reference still
% works; only the tones changed from the original green/blue accents.
\definecolor{kugreen}{HTML}{1A1A1A}   % near-black — section headings & title
\definecolor{kudark}{HTML}{333333}    % dark gray — subsection/subsubsection titles
\definecolor{headbg}{HTML}{EAEAEA}    % light gray background for table header rows
\definecolor{linkblue}{HTML}{404040}  % medium-dark gray — links / citations

% ---- Section heading styling -------------------------------------------------
% 'nobottomtitles*' stops a heading from being printed near the bottom of a
% page: if less than \bottomtitlespace of room is left, the heading (and its
% following text) moves to the next page. This is titlesec's built-in, robust
% replacement for the earlier \needspace hack — the latter injected vertical
% glue that clashed with longtable page-splitting ("Infinite glue shrinkage").
\usepackage[nobottomtitles*]{titlesec}
\titleformat{\section}
  {\color{kugreen}\Large\bfseries}{\thesection}{0.6em}{}
\titleformat{\subsection}
  {\color{kudark}\large\bfseries}{\thesubsection}{0.5em}{}
\titleformat{\subsubsection}
  {\color{kudark}\normalsize\bfseries}{\thesubsubsection}{0.5em}{}
\titlespacing*{\section}{0pt}{1.4\baselineskip}{0.5\baselineskip}

% Minimum room a heading needs at the bottom of a page; if only ~1–3 lines are
% left (less than this), the heading is bumped to the next page instead of being
% stranded above its table/text.
\renewcommand{\bottomtitlespace}{4\baselineskip}

% ---- Hierarchical indentation by heading level ------------------------------
% Each heading level shifts BOTH the heading and the body text under it further
% to the right: section = 0, subsection (x.y) = 1 step, subsubsection (x.y.z) =
% 2 steps. We do this by having the sectioning commands set \leftskip (a global
% left margin that carries through every following line until the next heading
% resets it). \par first, so the *previous* paragraph is not re-indented.
% \leftskip indents plain paragraphs; \@totalleftmargin makes list environments
% (itemize/enumerate) start from the same indent instead of the page margin.
\makeatletter
\newlength{\lvlindent}
\setlength{\lvlindent}{1.8em}          % one indent step
\let\ulmsSecIndent\section
\renewcommand{\section}{\par\global\leftskip=0pt\global\@totalleftmargin=0pt\relax\ulmsSecIndent}
\let\ulmsSubsecIndent\subsection
\renewcommand{\subsection}{\par\global\leftskip=\lvlindent\global\@totalleftmargin=\lvlindent\relax\ulmsSubsecIndent}
\let\ulmsSubsubsecIndent\subsubsection
\renewcommand{\subsubsection}{\par\global\leftskip=2\lvlindent\global\@totalleftmargin=2\lvlindent\relax\ulmsSubsubsecIndent}
\makeatother

% Indent the headings themselves to match (titlesec ignores \leftskip for the
% heading line, so its left margin is set here via the first \titlespacing arg).
\titlespacing*{\subsection}   {\lvlindent} {1.0\baselineskip}{0.4\baselineskip}
\titlespacing*{\subsubsection}{2\lvlindent}{0.8\baselineskip}{0.3\baselineskip}

% ---- Lists (tighter, custom bullets) ----------------------------------------
\usepackage{enumitem}
\setlist[itemize]{leftmargin=1.6em, itemsep=2pt, topsep=3pt}
\setlist[enumerate]{leftmargin=1.9em, itemsep=2pt, topsep=3pt}

% ---- Tables ------------------------------------------------------------------
\usepackage{array}
\usepackage{booktabs}         % \toprule / \midrule / \bottomrule
\usepackage{tabularx}         % auto-width columns that wrap long Thai text
\usepackage{longtable}        % tables that can break across pages
\usepackage{multirow}         % merged cells for the team-structure table

% Left-aligned wrapping column types:
%   L{width} — fixed-width ragged-right paragraph column
%   Y        — flexible ragged-right X column (for tabularx)
\newcolumntype{L}[1]{>{\raggedright\arraybackslash}p{#1}}
\newcolumntype{Y}{>{\raggedright\arraybackslash}X}
\renewcommand{\arraystretch}{1.3}

% Helpers for a styled header row inside tables.
%   \hrow  — starts a header row and shades it (must be FIRST in the row;
%            colortbl allows only one \rowcolor per row).
%   \thead — bold header-cell text.
\newcommand{\hrow}{\rowcolor{headbg}}
\newcommand{\thead}[1]{\textbf{#1}}

% \begin{keeptable} … \end{keeptable} — keeps a table and the bold label above
% it on the SAME page (used by the data dictionary in §5). A tabularx cannot
% split across pages, so LaTeX is free to break between the label and its
% table, stranding the label alone at the bottom of a page; wrapping the pair
% in a minipage makes them one unbreakable block.
%   * \leftskip is zeroed inside: the minipage box is already placed at the
%     current heading indent, so leaving it set would indent the content twice
%     and push the table past the right margin.
%   * \parskip is restored because minipage's \@parboxrestore zeroes it, which
%     would glue the table straight onto the label line.
% Tables inside must use \linewidth (= the minipage width), not \textwidth.
\newenvironment{keeptable}
  {\par\medskip\noindent
   \begin{minipage}{\dimexpr\textwidth-\leftskip\relax}%
   \leftskip=0pt\relax
   \setlength{\parskip}{0.5\baselineskip plus 2pt}}
  {\end{minipage}\par\medskip}

% ---- Table of contents styling ----------------------------------------------
% By default LaTeX gives dotted leaders to subsections but NOT to sections
% (they print bold with just a right-aligned page number). tocloft lets us turn
% on dots for the section level too, so every TOC line runs "….." to its page.
\usepackage{tocloft}
\renewcommand{\cftsecdotsep}{\cftdotsep}   % dotted leaders for section entries
\renewcommand{\cftsecleader}{\cftdotfill{\cftsecdotsep}}
\renewcommand{\cftsecpagefont}{\normalfont}  % keep page numbers un-bold

% ---- Figures -----------------------------------------------------------------
\usepackage{graphicx}
\graphicspath{{figures/}}

% Landscape (rotated) full-page floats — needed by the ER diagram in §5, which
% is far wider than it is tall and would be unreadable scaled to \textwidth.
\usepackage{rotating}

% ---- Flowcharts / diagrams (TikZ) -------------------------------------------
% Used for the borrow–return workflow diagram (§5). Libraries: geometric shapes
% for decision diamonds, arrows.meta for arrowheads, positioning for relative
% node placement, calc/fit/backgrounds for routing loops and grouping.
\usepackage{tikz}
\usetikzlibrary{shapes.geometric, arrows.meta, positioning, calc, fit, backgrounds}

% ---- Links & bookmarks -------------------------------------------------------
\usepackage{hyperref}
\hypersetup{
  colorlinks=true, linkcolor=black, citecolor=linkblue, urlcolor=linkblue,
  pdflang={th},
}

% Thai labels for auto-generated headings.
\renewcommand{\contentsname}{สารบัญ (Table of Contents)}
\renewcommand{\tablename}{ตาราง}
\renewcommand{\figurename}{รูปที่}

